\documentclass[11pt,a4paper]{article}
\usepackage[utf8]{inputenc}
\usepackage[T1]{fontenc}
\usepackage[french]{babel}
\frenchbsetup{StandardLists=true}
\usepackage[left=2cm,right=2cm,top=2cm,bottom=2cm]{geometry}
\usepackage[dvipsnames]{xcolor}
\usepackage{fouriernc}
\usepackage{amsmath,amssymb,amsfonts}
\usepackage{array,tabularx,booktabs,multirow}
\usepackage{colortbl}
\usepackage{graphicx}
\usepackage{mdframed}
\usepackage{enumitem}
\usepackage{fancyhdr}
\usepackage{chemformula}
\usepackage{awesomebox}
\setlength{\aweboxvskip}{0mm}
\usepackage{tcolorbox}
\tcbuselibrary{skins}
\usepackage{fontawesome5}

\definecolor{exoheader}{HTML}{1E5AA0}
\definecolor{exolight}{HTML}{DCE8FF}
\definecolor{warnorange}{RGB}{220,100,0}
\definecolor{problemebg}{HTML}{FFFDE7}

\renewcommand{\arraystretch}{1.8}
\setlength{\extrarowheight}{1mm}

\pagestyle{fancy}
\renewcommand\headrulewidth{1pt}
\fancyhead[L]{Académie de Besançon}
\fancyhead[C]{\textbf{TP --- Dosage de l'ibuprofène}}
\fancyhead[R]{Terminale / BTS Chimie}
\fancyfoot[C]{page \thepage}
\fancyfoot[R]{2025 -- 2026}
\fancyfoot[L]{Alexis RUIZ}

\newcommand{\sectitle}[1]{%
  \vspace{6pt}
  \noindent\colorbox{exoheader}{\parbox{\dimexpr\linewidth-2\fboxsep}%
    {\color{white}\bfseries\large\strut\quad #1}}
  \vspace{4pt}
}
\newcommand{\subsectitle}[1]{%
  \vspace{4pt}
  \noindent\colorbox{exolight}{\parbox{\dimexpr\linewidth-2\fboxsep}%
    {\color{exoheader}\bfseries\strut\quad #1}}
  \vspace{2pt}
}
\newcounter{question}
\newcommand{\question}{%
  \stepcounter{question}%
  \noindent\textcolor{exoheader}{\textbf{Question \thequestion.}}\quad
}

\begin{document}

\noindent\textbf{NOM :}\hfill\textbf{Prénom :}\hfill\textbf{Classe :}\hfill\phantom{x}\\[0.2cm]
\hrule\vspace{0.2cm}

\begin{center}
  {\large\bfseries\textcolor{exoheader}{Travaux Pratiques de Chimie}}\\[4pt]
  {\LARGE\bfseries Dosage de l'ibuprofène par pH-métrie (ExAO)}\\[2pt]
  {\small Académie de Besançon --- Belfort}
\end{center}
\hrule\vspace{6pt}

\begin{tcolorbox}[enhanced, colback=problemebg, colframe=exoheader, boxrule=1.5pt,
  title={\faFlask\quad Problématique},
  fonttitle=\bfseries\color{white},
  attach boxed title to top left={yshift=-2mm, xshift=4mm},
  boxed title style={colback=exoheader, colframe=exoheader}]
\textit{L'étiquette d'un comprimé d'ibuprofène indique 400 mg de principe actif.
\textbf{L'indication est-elle correcte ?}
On se propose de le vérifier par un dosage pH-métrique assisté par ordinateur (ExAO).}
\end{tcolorbox}

\vspace{6pt}

% ══════════════════════════════════════════════════════════════════════════════
\sectitle{Partie I --- Objectifs et principe}
% ══════════════════════════════════════════════════════════════════════════════

\begin{mdframed}[backgroundcolor=exolight, linecolor=exoheader, linewidth=1.2pt]
L'\textbf{ibuprofène} (\ch{C13H18O2}, $M = 206{,}3$~g/mol) est un acide faible
(anti-inflammatoire non stéroïdien). On le dose avec une solution de soude
\ch{NaOH} de concentration $C_b = 0{,}10$~mol/L.
\end{mdframed}

\vspace{4pt}
\question Quelle est la fonction chimique de l'ibuprofène ? Justifier. \dotfill

\vspace{0.3cm}\noindent\dotfill

\vspace{4pt}
\question Écrire l'équation de la réaction de titrage entre l'ibuprofène (\ch{AH}) et les ions hydroxyde. \dotfill

\vspace{0.3cm}\noindent\dotfill

% ══════════════════════════════════════════════════════════════════════════════
\sectitle{Partie II --- Protocole expérimental}
% ══════════════════════════════════════════════════════════════════════════════

\subsectitle{1. Préparation de la burette}

\begin{enumerate}[leftmargin=1.5cm, label=\textcolor{exoheader}{\textbf{\arabic*.}}]
  \item Placer le bécher étiqueté \og récupération \fg{} sous la burette.
  \item Vider la burette, puis la rincer avec la solution de soude.
  \item Remplir la burette avec la soude ($C_b = 0{,}10$~mol/L) et ajuster au zéro.
\end{enumerate}

\vspace{4pt}
\subsectitle{2. Dissolution du comprimé}

\begin{enumerate}[leftmargin=1.5cm, label=\textcolor{exoheader}{\textbf{\arabic*.}}]
  \item Peser précisément 1 comprimé d'ibuprofène. Noter la masse $m_{\text{comprimé}}$.
  \item Broyer finement le comprimé dans un mortier.
  \item Transvaser la poudre dans un bécher de 250~mL, ajouter 20~mL d'éthanol à 95\%.
  \item Agiter jusqu'à dissolution partielle, puis ajouter 80~mL d'eau distillée.
  \item Filtrer si nécessaire pour éliminer les excipients insolubles.
  \item Transvaser quantitativement dans une fiole jaugée de 100~mL ; compléter à la jauge.
        La solution obtenue est notée $S$ (concentration $C_a$ inconnue).
\end{enumerate}

\vspace{4pt}
\subsectitle{3. Prise d'essai et dosage ExAO}

\begin{enumerate}[leftmargin=1.5cm, label=\textcolor{exoheader}{\textbf{\arabic*.}}]
  \item Prélever $V_0 = 10{,}0$~mL de solution $S$ à la pipette jaugée et le verser dans un bécher.
  \item Ajouter 20~mL d'eau distillée et le barreau aimanté.
  \item Connecter la sonde pH-métrique à l'interface ExAO.
  \item Paramétrer le système : axe x = volume de soude versé (mL), axe y = pH.
  \item Effectuer le titrage en ajoutant 0,5~mL de soude par 0,5~mL (resserrer autour du saut de pH).
  \item Mettre fin à l'acquisition après la dernière mesure.
\end{enumerate}

\begin{mdframed}[backgroundcolor=orange!15, linecolor=warnorange, linewidth=2pt]
\textbf{\textcolor{warnorange}{\faExclamationTriangle\quad ATTENTION}} ---
L'éthanol est inflammable. Éviter toute flamme lors de la dissolution.
\end{mdframed}

% ══════════════════════════════════════════════════════════════════════════════
\newpage
\sectitle{Partie III --- Exploitation des résultats}
% ══════════════════════════════════════════════════════════════════════════════

\subsectitle{1. Lecture graphique}

\question Déterminer le volume équivalent $V_E$ (méthode des tangentes ou dérivée).

\vspace{0.3cm}
\begin{center}
\begin{tabularx}{0.7\linewidth}{|>{\centering\arraybackslash\columncolor{exolight}}X
                                |>{\centering\arraybackslash}X|}
\hline
\rowcolor{exoheader}
\color{white}\textbf{Grandeur} & \color{white}\textbf{Valeur} \\
\hline
Volume équivalent $V_E$ (mL) & \dotfill \\
\hline
pH à l'équivalence & \dotfill \\
\hline
\end{tabularx}
\end{center}

\vspace{4pt}
\subsectitle{2. Calcul de la concentration en ibuprofène}

\question En utilisant la relation à l'équivalence $C_b \times V_E = C_a \times V_0$,
calculer $C_a$ (en mol/L).

\vspace{2cm}

\noindent$C_a = $ \dotfill~mol/L

\vspace{4pt}
\subsectitle{3. Masse d'ibuprofène par comprimé}

\question Calculer la masse d'ibuprofène titrée dans les 10~mL prélevés, puis
en déduire la masse $m$ dans le comprimé entier.

\[
  m = C_a \times V_{\text{fiole}} \times M(\text{ibuprofène})
\]

\vspace{2cm}

\noindent$m = $ \dotfill~mg

\vspace{4pt}
\question Comparer avec la valeur théorique (400~mg). Calculer l'écart relatif.
Un écart $< 5\%$ est acceptable.

\vspace{0.3cm}\noindent\dotfill

\vspace{0.3cm}\noindent\dotfill

\vspace{4pt}
\subsectitle{4. Fin de manipulation}

\begin{enumerate}[leftmargin=1.5cm, label=\textcolor{exoheader}{\textbf{\arabic*.}}]
  \item Retirer la sonde pH-métrique, la rincer à l'eau distillée et la ranger dans sa solution de stockage.
  \item Vider la burette, la rincer à l'eau distillée et la remplir avec de l'eau distillée.
  \item Vider les béchers dans le récipient \og déchets \fg{}.
\end{enumerate}

% ══════════════════════════════════════════════════════════════════════════════
\newpage
\sectitle{Annexe --- Données}
% ══════════════════════════════════════════════════════════════════════════════

\vspace{4pt}
\begin{mdframed}[backgroundcolor=exolight, linecolor=exoheader, linewidth=1.2pt]
\begin{itemize}
  \item Ibuprofène : $M = 206{,}3$~g/mol ; formule brute \ch{C13H18O2} ; acide faible ($pK_a = 4{,}9$)
  \item Solution de soude : $C_b = 0{,}10$~mol/L
  \item Volume de la prise d'essai : $V_0 = 10{,}0$~mL
  \item Volume de la fiole jaugée : $V_{\text{fiole}} = 100$~mL
  \item Masse molaire de NaOH : 40~g/mol
\end{itemize}
\end{mdframed}

\end{document}
